\subsection{独立赋值}

定义 1. 设 A 与 A' 是同一域 K 的两个赋值环。若由 A 与 A' 生成的环是 K ，则称 A 与 A' 独立。K 上的两个赋值，若其环独立，则称为独立的，否则称为相依的。

K 上的非真赋值与 K 上的每个赋值独立。K 上的两个高为 1 的赋值独立，当且仅当它们不等价（ \(\S\) 4 第 5 节命题 6 (c)）。

定理 1 （赋值的逼近定理）. 设 \(v_{i}\) ( \(1 \leqslant i \leqslant n\) ) 是域 \(\mathbf{K}\) 上两两独立的赋值，\(\Gamma_{i}\) 是 \(v_{i}\) 的阶群。设 \(a_{i} \in \mathbf{K}\) 且 \(\alpha_{i} \in \Gamma_{i}\) ( \(1 \leqslant i \leqslant n\) )。那么存在 \(x \in \mathbf{K}\) 使得 \(v_{i}(x - a_{i}) \geqslant \alpha_{i}\) 对一切 \(i\) 成立。

若 \(v_{i}\) 是非真的，则 \(\alpha_{i} = 0\) ，且关系 \(v_{i}(x - a_{i}) \geqslant \alpha_{i}\) 对一切 \(x \in \mathbf{K}\) 成立。因此可以设诸 \(v_{i}\) 都不是非真的。

设 \(\mathbf{A}_i\) 是 \(v_i\) 的环，\(\mathbf{B} = \bigcap_{i=1} \mathbf{A}_i\) ，\(\mathfrak{p}_i = \mathfrak{m}(\mathbf{A}_i) \cap \mathbf{B}\) 。由第 1 节命题 1，诸 \(a_i\) 可写成 \(a_i = b_i / s\) ( \(b_i \in \mathbf{B}\) ，\(s \in \mathbf{B} - \{0\}\) )；若记 \(x = y / s\) 且 \(\alpha_i' = \alpha_i + v_i(s)\) ，则 \(v_i(y - b_i) \geqslant \alpha_i'\) 。这表明可以设 \(a_i \in \mathbf{B}\) 对一切 \(i\) 成立；也可以设 \(\alpha_i > 0\) 对一切 \(i\) 成立。设 \(\mathfrak{v}_i\) 是由这样的 \(z \in K\) 组成的集合：\(v_{i}(z) \geqslant \alpha_{i}\) ；记 \(q_{i} = v_{i} \cap B\) 。对 \(x \in B\) ，\(v_{i}(x - a_{i}) \geqslant \alpha_{i}\) 等价于 \(x \equiv a_{i}(q_{i})\) 。因此只需证明典范同态

\(\mathbf{B} \to \prod_{i=1}^{n} \left( \mathbf{B} / q_i \right)\) 是满射，即 \(q_i + q_j = \mathbf{B}\) 对 \(i \neq j\) 成立（第二章 § 1 第 2 节命题 5）。由于 \(\mathbf{B}\) 的极大理想就是诸 \(p_i\) （命题 2），为此只需证明 \(q_i \not\subset p_j\) 对 \(i \neq j\) 成立。

设存在 \(i, j\) 使得 \(\mathfrak{q}_i \subset \mathfrak{p}_j\) 且 \(i \neq j\) 。我们马上会看到 \(\mathfrak{q}_i\) 的根是 B 的素理想 \(\mathfrak{p}\) 。那么 \(\mathfrak{p} \subset \mathfrak{p}_j\) ，且也有 \(\mathfrak{p} \subset \mathfrak{p}_i\) ，因为 \(\alpha_i > 0\) 从而 \(\mathfrak{q}_i \subset \mathfrak{p}_i\) 。因此 \(\mathrm{A}_j = \mathrm{B}_{\mathfrak{p}_j} \subset \mathrm{B}_{\mathfrak{p}}\) （第 1 节命题 1），类似地 \(\mathrm{A}_i \subset \mathrm{B}_{\mathfrak{p}}\) 。而现在，由于 \(\mathfrak{v}_i \neq (0)\) 且 \(\mathfrak{v}_i = \mathrm{B}_{\mathfrak{p}_i} \mathfrak{q}_i\) （第二章 §2 第 4 节命题 10），有 \(\mathfrak{q}_i \neq (0)\) ，从而 \(\mathfrak{p} \neq (0)\) 且 \(\mathrm{B}_{\mathfrak{p}} \neq \mathrm{K}\) 。这与 \(\mathrm{A}_i\) 和 \(\mathrm{A}_j\) 独立的假设矛盾。

剩下的要证明 p 是素的。而这由下面的引理得出：

引理 2. 设 A 是赋值环，b 是 A 的异于 A 的理想。那么 b 的根 r 是素理想。

设 \(xy \in \mathfrak{r}\) 。那么存在 \(n \geqslant 1\) 使得 \((xy)^n \in \mathfrak{b}\) 。设 \(v\) 是与 A 相伴的赋值。例如若 \(v(x) \geqslant v(y)\) ，则

\[
v (x ^ {2 n}) \geqslant v (x ^ {n} y ^ {n}),
\]

于是 \(x^{2n} \in \mathfrak{b}\) ，从而 \(x \in \mathfrak{r}\) 。

推论 1. 对每个由元素 \(\gamma_{i} \in \Gamma_{i}\) ( \(1 \leqslant i \leqslant n\) ) 组成的族，存在 \(x \in \mathbf{K}\) 使得 \(v_{i}(x) = \gamma_{i}\) ( \(1 \leqslant i \leqslant n\) )。

可以设 \(\mathbf{A}_i \neq \mathbf{K}\) 对一切 \(i\) 成立。那么对一切 \(i\) 存在 \(a_i \in \mathbf{K}\) 使得 \(v_i(a_i) = \gamma_i\) ，且存在 \(\alpha_i \in \Gamma_i\) 使得 \(\gamma_i < \alpha_i\) 。对这些元素 \(a_i\) 应用定理 1：存在 \(x \in \mathbf{K}\) 使得 \(v_i(x - a_i) > v_i(a_i)\) ；于是由 \(x = a_i + (x - a_i)\) 得 \(v_i(x) = v_i(a_i) = \gamma_i\) （ \(\S 3\) 第 1 节命题 1）。

推论 2. 设 \(\mathcal{T}_i\) 是 \(\mathbf{K}\) 上由 \(v_i\) 定义的拓扑；给 \(\mathbf{K}^n\) 以诸 \(\mathcal{T}_i\) 的乘积拓扑。若诸 \(v_i\) 不是非真的，则 \(\mathbf{K}^n\) 的对角线在 \(\mathbf{K}^n\) 中稠密。

命题 3. 设 \(v\) 与 \(v'\) 是同一域 \(K\) 上两个不是非真的赋值。\(v\) 与 \(v'\) 在 \(K\) 上定义同一个拓扑，当且仅当它们相依。

设拓扑 \(\mathcal{T}_v\) 与 \(\mathcal{T}_{v'}\) （分别由 \(v\) 与 \(v'\) 定义）相同。由于 \(\mathcal{T}_v\) 是 Hausdorff 的，\(\mathbf{K}^2\) 的对角线是闭的，因而 \(v\) 与 \(v'\) 相依（定理 1 的推论 2）。

反之，设 \(v\) 与 \(v'\) 相依。那么它们的环 A 与 A' 含于同一个异于 K 的环 A'' ，且 A'' 是某个赋值 \(v''\) 的环（§ 4 第 1 节命题 1）。只需证明拓扑 \(\mathcal{T}_{v''}\) 与 \(\mathcal{T}_v\) 相同。设 \(\Gamma\) 与 \(\Gamma''\) 是 \(v\) 与 \(v''\) 的阶群。存在一个递增同态 \(\lambda\) （由 \(\Gamma\) 到 \(\Gamma''\) 上）使得 \(v'' = \lambda \circ v\) （§ 4 第 3 节）。若 \(\alpha'' \in \Gamma''\) ，取 \(\alpha \in \overline{\lambda}^{1}(\alpha'')\) ；条件 \(v(x) \geqslant \alpha\) 蕴含 \(v''(x) \geqslant \alpha''\) 。设 \(\beta \in \Gamma\) 且 \(\beta'' = \lambda(\beta)\) ；条件 \(v(x) \leqslant \beta\) 蕴含 \(v''(x) \leqslant \beta''\) ，因而条件 \(v''(x) > \beta''\) 蕴含 \(v(x) > \beta\) 。由于 \(v\) 与 \(v''\) 不是非真的，所论的不等式分别定义了 \(\mathcal{T}_v\) 与 \(\mathcal{T}_{v''}\) 的 0 的基本邻域系。于是 \(\mathcal{T}_v = \mathcal{T}_{v''}\) ，这就完成了证明。

注

\begin{enumerate}
\def\labelenumi{(\arabic{enumi})}
\item
  命题 3 表明关系 ``v 与 v' 相依'' 是一个等价关系。
\item
  考虑到高为 1 的赋值与超度量绝对值之间的关系（§ 6 第 2 节），在高为 1 的赋值的情形，命题 3 也可由等价绝对值的刻画（《一般拓扑学》第九章 § 3 第 2 节命题 5）得出。
\end{enumerate}

命题 4. 设 \(v_{1}, \ldots, v_{n} (n \geqslant 2)\) 是同一域 K 上两两相依的赋值。那么环 \(A_{1}, \ldots, A_{n}\) （即 \(v_{1}, \ldots, v_{n}\) 的环）生成 K 的一个异于 K 的子环。

对 \(n = 2\) ，命题 4 由定义 1 得出。设它对 \(n - 1\) 个赋值成立。那么存在 K 的异于 K 且包含 \(\mathbf{A}_1, \ldots, \mathbf{A}_{n-1}\) 的子环 A ；也存在一个子环 \(\mathbf{B} \neq \mathbf{K}\) ，它包含 \(\mathbf{A}_{n-1}\) 与 \(\mathbf{A}_n\) 。由于 A 与 B 都包含 \(\mathbf{A}_{n-1}\) ，它们按包含关系是可比较的（§ 4 第 1 节命题 1 的推论）。因此两者中较大者包含所有 \(\mathbf{A}_i\) 。

